\section{Partie théorique, Ondes et interférences} \label{sec:theorie}

Soit une onde indéformable :
\begin{equation}
    \begin{center}
    $p(x,t) = Â \cdot sin(\omega t - kx)$ 
    \caption{eq:onde}
    \end{center}
\end{equation}

qui dépend d'une composante spaciale (x) et temporelle (t). Cette onde est définie par k le nombre d'onde et $\omega$ la pulsation tels que :
\begin{equation}
    \begin{center}
        k = \frac{2\pi}{\lambda} [...] et \omega = 2\pi f [rad/s]
    \caption{eq:param}
    \end{center}
\end{equation}

p(x,t) est donc définie plus spécifiquement par une longueur d'onde $\lambda$ [m], une fréquence f [Hz] et une amplitude Â.
\\\\
Nous nous interessons a deux cas specifiques. En un premier temps aux interferences d'ondes accoustiques, puis aux interferences d'ondes de surfaces circulaires. 

\subsection{Interférences d'ondes indéformables}
Soit deux ondes indéformables p1(x,t) et p2(x,t). L'onde stationnaire résultant de la superposition de p1 et p2, soit par la somme des deux tel que :
\begin{equation}
    p_m(t) = p1 + p2 = 2 Â cos(k \frac{d_2 - d_1}{2} sin(\omega t - k \frac{d_1 + d_2}{2})
    \caption{eq:station}
\end{equation}

Pour deux ondes cohérentes de même amplitude, l'interférence est donc constructive si la différence de chemin \(\Delta x\) est un multiple entier de \(\lambda\) :
\begin{equation}
  \Delta x = n \lambda \quad \text{avec } n \in \mathbb{Z}
  \] 
\end{equation}
et destructive si la différence de chemin \(\Delta x\) est un multiple demi-entier de \(\lambda\) :
\begin{equation}
    \Delta x = \left(n + \frac{1}{2}\right) \lambda \quad \text{avec } n \in \mathbb{Z}
    \caption{eq:foo}
\end{equation}

Nous pouvons observer ce phenomène graphiquement sur Fig. \ref{fig:ohwaw} :

\begin{figure}[H]
\begin{center}
\includegraphics[scale=0.4]{ohwaw.png}
\end{center}
\label{fig:ohwaw}
\caption{Illustration d'ondes constructives et destructives}
\end{figure}

% Include une images pour les deux ?


\subsection{Ondes acoustiques}

Les ondes acoustiques suivent la forme générale d'une onde indéformable présente sur Eq. \ref{eq:onde}
\\\\
A amplitude constante \(A\) ne dépend pas de la position \(x\). Les nœuds correspondent à l'amplitude minimale et sont espacés de \(\frac{\lambda}{2}\). Ce comportement est propre des ondes stationnaires dans un milieu unidimensionnel. Par définition Ref. \cite{protocole}, la célérité est définite tel que le rapport $\frac{\omega}{k}$, soit par l

\subsection{Ondes de surface circulaires}

Les ondes de surface circulaires partagent la même base théorique que les ondes acoustiques, mais leur propagation est influencée par la géométrie du milieu.

\paragraph{Équation et amplitude}
L'onde s'écrit :
\[
h(r,t) = \frac{A_0}{\sqrt{r}} \cdot \sin(\omega t - kr)
\]
où \(A_0\) est l'amplitude initiale et \(r\) la distance radiale. **L'amplitude décroît avec \(\frac{1}{\sqrt{r}}\)** en raison de la conservation de l'énergie sur une surface circulaire.

\paragraph{Célérité}
La vitesse de propagation \(v\) dépend de la profondeur \(d\) et de la longueur d'onde \(\lambda\) :
- En eau profonde (\(d \gg \lambda\)) :
  \[
  v = \sqrt{\frac{g \lambda}{2\pi}}
  \]
- En eau peu profonde (\(d \ll \lambda\)) :
  \[
  v = \sqrt{g d}
  \]
où \(g\) est l'accélération gravitationnelle.

\paragraph{Particularités}
- **Amplitude variable** : Contrairement aux ondes acoustiques, \(A\) dépend de la distance \(r\).
- **Interférences** : Les franges (zones d'amplitude maximale/minimale) forment des motifs circulaires, avec une distance entre franges brillantes égale à \(\lambda\)



















%\subsection{Ondes accoustiques}
%
%Une onde indeformable de pression accoustique peut s'ecrire par...
%
%Pour les ondes acoustiques, la célérité \textbf{c} d'une onde accoustique dans un milieu est donné par :
%\[
%    c = \sqrt{\gamma \frac{R M}{T}} 
%\]
%Avec gamma ..., R ... et M .....
%\\\\
%De manière plus générale, la vitesse de propagation \(v\) d'une onde est donnée par la relation :
%\[
%c = \lambda \cdot f
%\]
%
%
%\subsection{Ondes de surface circulaires}
%
%De même que pour les ondes de pression accoustiques, une onde de surface circulaire peut s'ecrire tel que : ...
%\\\\
%A la seule différence que sa nature géométrique influence l'amplitude de l'onde tel que pour : ... 
%

